\section*{Hoofdstuk \ref{ch:b2:viscous-flows} --- Viskeuze stromingen}

\begin{solution}{exo:b2:viscous-flows:1}
$\tau = \eta U/h = 0.25 \times 2/2 \times 10^{-4} = \qty{2.5}{kPa}$;
$F = \qty{1.25}{kN}$; $P = Fv = \qty{2.5}{kW}$. Droge wrijving:
$0.3 \times 98 = \qty{29}{N}$ --- hier verzet de dunne, sterk afgeschoven
film zich veel meer; de viskeuze weerstand groeit als $U/h$, zodat smering
loont bij lagere snelheden en dikkere films (en zij loopt nooit vast).
\end{solution}

\begin{solution}{exo:b2:viscous-flows:2}
(a) $2 \times 10^{-6} \times 3 \times 10^{-5}/10^{-6} = 6 \times 10^{-5}$,
kruipend. (b) $1.8 \times 1.5/10^{-6} = 2.7 \times 10^6$, turbulent. (c)
$4 \times 30/1.5 \times 10^{-5}
= 8 \times 10^6$, turbulent. (d)
$0.025 \times 0.33/3.3 \times 10^{-6} = 2500$, net laminair. (e)
$900 \times 0.5 \times 0.01/0.1 = 45$, laminair.
\end{solution}

\begin{solution}{exo:b2:viscous-flows:3}
$D_V = \pi R^4\Delta P/8\eta L = \pi(2 \times 10^{-4})^4 \times 2 \times 10^4/(8 \times 10^{-3}
\times 0.03) = \qty{0.42}{mL/s}$;
\qty{12}{s} voor \qty{5}{mL}.
$v_{\max} = \Delta PR^2/
4\eta L = \qty{6.7}{m/s}$,
$\langle v\rangle = \qty{3.3}{m/s}$,
$\mathrm{Re} = 2R\langle v
\rangle/\nu = 1300$: laminair.
$\tau_w = \Delta PR/2L = \qty{67}{Pa}$. Bij \qty{0.8}{mm}: zestien keer het
debiet, \qty{6.7}{mL/s} --- maar $\mathrm{Re} \approx 10^4$: de stroming
wordt turbulent en de winst is kleiner.
\end{solution}

\begin{solution}{exo:b2:viscous-flows:4}
$v_\infty = 2R^2\Delta\rho\,g/9\eta$. (a) \qty{1.3}{cm/s},
$\mathrm{Re} = 0.02$, \qty{75}{s} per meter. (b) \qty{0.24}{mm/s},
$\mathrm{Re} = 3 \times 10^{-5}$, \qty{70}{min} per meter. (c)
\qty{1.2}{m/s} maar $\mathrm{Re} = 16$: Stokes overschat (de werkelijke
snelheid is ongeveer \qty{0.7}{m/s}), ongeveer een seconde per meter.
\end{solution}

\begin{solution}{exo:b2:viscous-flows:5}
(a) $\langle v\rangle = 0.1/0.0707 = \qty{1.4}{m/s}$,
$\mathrm{Re} = 4.2 \times 10^5$. (b) $8\eta LD_V/\pi R^4 = \qty{500}{Pa}$.
(c) $\lambda = 0.3/25.5 = 0.0118$,
$\Delta P = 0.0118
\times 3333 \times 1000 = \qty{39}{kPa}$: $80$ keer meer;
$P = \Delta PD_V = \qty{3.9}{kW}$. (d) Dezelfde snelheid,
$\mathrm{Re} = 3 \times 10^5$, $\lambda = 0.0128$,
$\Delta P = 0.0128 \times 4717
\times 1000 = \qty{60}{kPa}$, \qty{6}{kW}:
meer wand per eenheid debiet, meer verlies.
\end{solution}

\begin{solution}{exo:b2:viscous-flows:6}
(a)
$v = 2 \times 25 \times 10^{-12} \times 50 \times 9.81/(9 \times 10^{-3}) = \qty{2.7}{\micro m/s}$;
\qty{5}{cm} in $1.8 \times 10^4$ s, vijf uur. (b) \qty{2.7}{mm/s},
\qty{18}{s}. (c)
$\Omega = \sqrt{1000g/r} = \qty{313}{rad/s} = \qty{3000}{rpm}$. (d)
$v \propto R^2$: grote deeltjes slaan het eerst neer; opeenvolgende
centrifugegangen met stijgend toerental scheiden de maten.
\end{solution}

\begin{solution}{exo:b2:viscous-flows:7}
(a) $F = 0.6 \times 144 \times 0.35 = \qty{30}{N}$, $+\qty{4}{N}$:
$P = 34 \times 12 =
\qty{410}{W}$. (b) $47 + 4 = \qty{51}{N}$, \qty{770}{W}.
(c) Relatieve snelheid \qty{17}{m/s}: $61 + 4 = \qty{65}{N}$, maal de
grondsnelheid: \qty{780}{W}. (d) Luchtweerstand $\qty{21}{N}$: \qty{300}{W},
een besparing van \qty{110}{W}. (e)
$mg\sin\alpha = \qty{39}{N} =
4 + 0.21v^2$: $v = \qty{13}{m/s}$,
\qty{47}{km/h}.
\end{solution}

\begin{solution}{exo:b2:viscous-flows:8}
(a) Golfbal: $\tfrac12\rho C_xS = \qty{2.2e-4}{kg/m}$,
$v_\infty = \sqrt{0.45/2.2 \times 10^{-4}}
= \qty{45}{m/s}$; parachutist
$\sqrt{785/0.42} = \qty{43}{m/s}$; met parachute
$\sqrt{785/15} =
\qty{7.2}{m/s}$. (b) $m\dot v = mg - kv^2$, d.w.z.\
$\dot v = g(1 - v^2/v_\infty^2)$, met als oplossing vanuit rust
$v_\infty\tanh(gt/v_\infty)$. (c) $\tanh x = 0.9$ bij $x = 1.47$:
$t = 1.47 \times 43/9.81 = \qty{6.5}{s}$;
$z = (v_\infty^2/g)\ln\cosh x =
190 \times 0.83 = \qty{160}{m}$.
\end{solution}

\begin{solution}{exo:b2:viscous-flows:9}
(a) Afschuifsnelheid $\Omega R/e = 0.5/10^{-3} = \qty{500}{s^{-1}}$;
$\Gamma = \tau\cdot2\pi Rh
\cdot R$:
$\tau = 0.02/(2\pi \times 2.5 \times 10^{-3} \times 0.1) = \qty{12.7}{Pa}$.
(b) $\eta
= 12.7/500 = \qty{2.5e-2}{Pa.s}$. (c)
$\mathrm{Re} = \rho Ue/\eta \approx 1000 \times 0.5 \times
10^{-3}/0.025 = 20$:
laminair. (d) Een dunne spleet maakt de afschuifsnelheid uniform (de vlakke
benadering) en houdt de stroming laminair.
\end{solution}

\begin{solution}{exo:b2:viscous-flows:10}
$D_V = \qty{8.3e-5}{m^3/s}$. (a)
$R_h = 8\eta L/\pi R^4 = 8 \times 3.5 \times 10^{-3} \times
0.4/(\pi \times 2.4 \times 10^{-8}) = \qty{1.5e5}{Pa.s/m^3}$,
$\Delta P = \qty{12}{Pa}$. (b) Eén haarvat $\qty{3.5e16}{Pa.s/m^3}$;
$10^{10}$ parallel: \qty{3.5e6}{Pa.s/m^3}; $\Delta P = \qty{290}{Pa}$
(\qty{2}{mmHg}). (c) Eén arteriole
$8 \times 3.5 \times 10^{-3}
\times 2 \times 10^{-3}/(\pi \times 10^{-20}) = \qty{1.8e15}{Pa.s/m^3}$;
het geheel: \qty{1.8e7}{Pa.s/m^3}; $\Delta P = \qty{1.5}{kPa}$: de
arteriolen overheersen de twee andere (de drukval zit in de kleine vaten,
niet in de grote). (d) $P = 13000 \times 8.3 \times 10^{-5} = \qty{1.1}{W}$,
ongeveer $1\%$. (e) $(1/1.1)^4 = 0.68$: een derde van de arteriolaire
weerstand weg --- zo verlagen vaatverwijders de bloeddruk.
\end{solution}

\begin{solution}{exo:b2:viscous-flows:11}
(a)
$F = b\int_0^L\tau_w\,\dd x = 0.332\rho U^2b\sqrt{\nu/U}\int_0^L\dd x/\sqrt x =
0.664\rho U^2b\sqrt{\nu L/U} = 0.664\rho U^2bL/\sqrt{\mathrm{Re}_L}$.
(b) $\mathrm{Re}_L =
3 \times 10^5$ (laminair);
$\delta \approx 5\sqrt{\nu L/U} = \qty{9}{mm}$;
$F = 0.664 \times 1000
\times 0.09 \times 1/548 = \qty{0.11}{N}$ per vlak,
\qty{0.22}{N}. (c) $\mathrm{Re}_L = 10^7$: de laag is voorbij de eerste
centimeters turbulent en de laminaire formule faalt (de werkelijke weerstand
is enkele malen groter). (d) $\mathrm{Re}_x =
10^9$:
$\delta \approx 0.37 \times 100/63 = \qty{0.6}{m}$.
\end{solution}

\begin{solution}{exo:b2:viscous-flows:12}
(a) $0 = \rho g + \eta v''(y)$, $v(0) = 0$ (hechting), $v'(h) = 0$ (geen
spanning aan het vrije oppervlak). (b) Twee keer integreren:
$v = (\rho g/2\eta)(2hy - y^2)$; $q =
\int_0^hv\,\dd y = \rho gh^3/3\eta$.
(c)
$v_s = \rho gh^2/2\eta = 1200 \times 9.81 \times
4 \times 10^{-8}/1 = \qty{0.47}{mm/s}$;
$q = \qty{6.3e-8}{m^2/s}$; \qty{28}{cm} in tien minuten --- een newtoniaanse
verf van die dikte zou eraf lopen; echte verven worden dunner bij
afschuiving en zijn thixotroop, en zakken toch uit als zij te dik zijn,
omdat $v_s \propto h^2$. (d)
$1400 \times 9.81 \times 4 \times 10^{-6}/20 = \qty{2.7}{mm/s}$.
\end{solution}

\begin{solution}{pb:b2:viscous-flows:1}
\textbf{1.} $S = \qty{4.9}{cm^2}$,
$\langle v\rangle = 83/4.9 = \qty{0.17}{m/s}$;
$\mathrm{Re}
= 1060 \times 0.17 \times 0.025/3.5 \times 10^{-3} = 1300$:
laminair (de pulserende pieken halen $\sim 5000$).

\textbf{2.} $v_{\max} = \qty{0.34}{m/s}$;
$\tau_w = 4\eta\langle v\rangle/R = \qty{0.19}{Pa}$.

\textbf{3.}
$\mathrm{Re} = 1060 \times 3 \times 10^{-4} \times 8 \times 10^{-6}/3.5 \times 10^{-3} =
7 \times 10^{-4}$;
$\Delta P = 8\eta L\langle v\rangle/R^2 = 8 \times 3.5 \times 10^{-3} \times 10^{-3}
\times 3 \times 10^{-4}/1.6 \times 10^{-11} = \qty{525}{Pa} = \qty{3.9}{mmHg}$.

\textbf{4.} $1/0.3 = \qty{3.3}{s}$: lang genoeg opdat zuurstof de enkele
micrometers naar het weefsel kan diffunderen
(\cref{ch:b2:particle-diffusion}).

\textbf{5.} Eén arteriole $8\eta L/\pi R^4 = \qty{1.8e15}{Pa.s/m^3}$; $10^8$
parallel: \qty{1.8e7}{Pa.s/m^3};
$\Delta P = 1.8 \times 10^7 \times 8.3 \times 10^{-5} =
\qty{1.5}{kPa} = \qty{11}{mmHg}$.

\textbf{6.}
$P = \Delta P\,D_V = 13000 \times 8.3 \times 10^{-5} = \qty{1.1}{W}$,
ongeveer $1\%$ van het vermogen van het lichaam.

\textbf{7.} $(1/0.6)^4 = 7.7$. In rust verwijden de vaten stroomafwaarts
zich en compenseren dat; bij inspanning stijgt de vraag vier- tot vijfvoudig
en kan de vernauwde slagader die niet leveren: angina pectoris.

\textbf{8.} Vijf keer de snelheid: $\mathrm{Re} \approx 6500$, turbulent ---
de geruisen die een stethoscoop tijdens inspanning hoort.

\textbf{9.} In de trage stroming met lage afschuiving in kleine aders en
haarvaten klonteren de rode bloedcellen samen en stijgt de schijnbare
\omterm{def:b2:viscous-flows:viscosity}{viscositeit} (in de snelle aorta juist andersom); alle schattingen van de
drukval in kleine vaten zijn dus ondergrenzen.

\textbf{10.} $\langle v\rangle = 3/1.13 = \qty{2.65}{m/s}$;
$\mathrm{Re} = 850 \times 2.65 \times
1.2/0.02 = 1.35 \times 10^5$:
turbulent.

\textbf{11.} $\lambda = 0.316/19.2 = 0.0165$;
$\Delta P/L = (0.0165/1.2) \times \tfrac12 \times
850 \times 7.0 = \qty{41}{Pa/m} = \qty{41}{kPa/km}$;
\qty{534}{bar} over de leiding.

\textbf{12.} $534/80 = 6.7$: minstens zeven stations.

\textbf{13.}
$P = \Delta P\,D_V = 5.34 \times 10^7 \times 3 = \qty{160}{MW}$;
\qty{53}{MJ/m^3}, $0.15\%$ van de energie van de olie.

\textbf{14.}
$\Delta P = 32\eta L\langle v\rangle/D^2 = 32 \times 0.02 \times 1.3 \times 10^6
\times 2.65/1.44 = \qty{15}{bar}$:
vijfendertig keer minder. Turbulentie is duur; enkele delen per miljoen aan
lange polymeren stellen haar uit en verlagen de val met tientallen
procenten.

\textbf{15.} $\mathrm{Re} = 1350$: laminair; Poiseuille:
$\Delta P = 32 \times 2 \times 1.3 \times
10^6 \times 2.65/1.44 = \qty{1500}{bar}$
--- onmogelijk; de olie moet warm worden gehouden, vandaar de isolatie en de
verwarming.

\textbf{16.} $1.3 \times 10^6/2.65 = \qty{4.9e5}{s} = \qty{5.7}{dagen}$.

\textbf{17.} $\tau_w = (\Delta P/L)D/4 = 41 \times 0.3 = \qty{12}{Pa}$;
$F = \tau_w\pi DL =
\qty{6e7}{N}$.

\textbf{18.}
$\Delta P \propto \lambda\langle v\rangle^2/D \propto D^{-4.75}$:
$1.2^{4.75} = 2.4$.

\textbf{19.}
$v = 2R^2\Delta\rho\,g/9\eta = 2 \times 2.5 \times 10^{-9} \times 1800 \times 9.81/
0.18 = \qty{0.49}{mm/s}$;
$\mathrm{Re} = 2 \times 10^{-3}$; \qty{41}{min} om \qty{1.2}{m} over te
steken. Turbulente schommelingen van enkele centimeters per seconde houden
de korrel zwevend zolang de olie loopt; in een stilgelegde leiding bezinkt
zij binnen het uur.

\textbf{20.}
$v = 2 \times 2.5 \times 10^{-11} \times 1000 \times 9.81/1.62 \times 10^{-4} =
\qty{3.0}{mm/s}$;
$\mathrm{Re} = 2 \times 10^{-3}$; \qty{1}{km} in $3.3 \times 10^5$ s, vier
dagen: elke opwaartse stroom boven enkele millimeters per seconde houdt de
druppel vast.

\textbf{21.} \qty{1.2}{m/s}, $\mathrm{Re} = 16$: Stokes overschat met
ongeveer een factor twee (in werkelijkheid \qty{0.7}{m/s}).

\textbf{22.} Stokes zou \qty{120}{m/s} en $\mathrm{Re} \sim 10^4$ geven:
absurd. Kwadratische weerstand,
$\tfrac43\pi R^3\rho_wg = \tfrac12\rho v^2\pi R^2C_x$:
$v = \sqrt{8R
\rho_wg/3\rho C_x} = \qty{6.6}{m/s}$ ($\mathrm{Re} = 900$,
consistent); \qty{150}{s} vanaf \qty{1}{km}.

\textbf{23.} \qty{1}{cm/s} $>$ \qty{3}{mm/s}: de druppel stijgt; de
motregendruppel heeft ongeveer \qty{1}{m/s} nodig, wat convectieve wolken
leveren.

\textbf{24.} $2\gamma/R$: \qty{29}{kPa}, \qty{1.4}{kPa}, \qty{144}{Pa};
dynamische drukken $\tfrac12\rho v^2$: $5 \times 10^{-6}$, $0.3$ en
\qty{26}{Pa}. De verhouding bereikt $0.2$ bij \qty{2}{mm}; bij \qty{5}{mm}
($v \approx \qty{9}{m/s}$, $\tfrac12\rho v^2 =
\qty{49}{Pa}$ tegen
\qty{58}{Pa}) plat de aerodynamische druk de druppel af en breekt hem: geen
regendruppel wordt groter dan ongeveer \qty{5}{mm}.

\textbf{25.} Aorta $10^3$ (laminair, dicht bij de overgang); haarvat
$10^{-3}$ (Poiseuille); arteriole $10^{-2}$ (Poiseuille); warme pijpleiding
$10^5$ (turbulent, $\lambda(\mathrm{Re})$); koude pijpleiding $10^3$
(Poiseuille); zand in olie $10^{-3}$ (Stokes); wolkendruppel $10^{-3}$
(Stokes); regendruppel $10^3$ (kwadratische weerstand, $C_x \approx 0.5$).
\end{solution}
